% -----------------------------
\section{MOTIVATION}
%------------------------------


\begin{comment}

A ocupação do ambiente marinho por estruturas offshore tem crescido de forma significativa nas últimas décadas, impulsionada principalmente pela exploração e produção de petróleo e gás. No Brasil, essa atividade concentra-se majoritariamente nas bacias sedimentares da margem continental sudeste e nordeste, com destaque para a Bacia de Campos, a Bacia de Santos, a Bacia do Espírito Santo e a Bacia Potiguar, sob operação predominante da Petrobras.

Essas regiões apresentam uma ampla variedade de condições oceanográficas e geotécnicas, abrangendo desde áreas offshore rasas, com profundidades inferiores a $100\,\mathrm{m}$, até ambientes de águas profundas e ultraprofundas, que ultrapassam $2{.}000\,\mathrm{m}$ de lâmina d’água, especialmente na província do pré-sal. A instalação dessas estruturas envolve etapas construtivas de alta energia, como a cravação de estacas (\textit{pile-driving}), a fixação de \textit{jackets}, a ancoragem de unidades flutuantes e a instalação de dutos submarinos, todas potencialmente associadas à geração de vibrações mecânicas e ruído subaquático.

%====================================================
% FIGURA mostrando os tipos de estruturas  offshore
%====================================================



Do ponto de vista ambiental, essas áreas coincidem com regiões de elevada biodiversidade marinha. A fauna local inclui diversas espécies de peixes demersais e recifais, comunidades bentônicas sensíveis a perturbações do sedimento e, de forma particularmente relevante, grandes mamíferos marinhos. Entre eles destacam-se a baleia-jubarte, que utiliza a costa brasileira como área reprodutiva durante seu ciclo migratório, a baleia-franca-austral, além de golfinhos costeiros e offshore, como o golfinho-nariz-de-garrafa e o boto-cinza.

Esses mamíferos dependem fortemente do ambiente acústico para comunicação, orientação espacial, alimentação e interação social. Assim, alterações no campo acústico marinho, mesmo que temporárias, podem provocar desde respostas comportamentais — como evitação de áreas e mudanças em rotas migratórias — até efeitos fisiológicos mais severos, a depender da intensidade, duração e conteúdo espectral do ruído gerado. Em ambientes offshore brasileiros, caracterizados por grandes profundidades e sedimentos predominantemente moles, a propagação de ondas acústicas e vibratórias pode ocorrer em escalas espaciais significativas, ampliando o alcance potencial desses impactos.

Nesse contexto, torna-se essencial o desenvolvimento de estudos que busquem compreender, modelar e prever a geração e a propagação de vibrações associadas à instalação de estruturas offshore, bem como suas interações com o meio marinho. Ainda que tais estudos atuem em uma escala aparentemente técnica ou restrita, eles constituem uma etapa fundamental para a formulação de estratégias de mitigação mais eficazes, capazes de reduzir impactos sobre espécies sensíveis e contribuir para uma convivência mais equilibrada entre o avanço tecnológico e a preservação ambiental.

Além das bacias offshore tradicionalmente exploradas no sudeste brasileiro, destaca-se a chamada Margem Equatorial brasileira \ref{fig:equatorial}, que se estende aproximadamente do Amapá ao Rio Grande do Norte e engloba bacias como Foz do Amazonas, Pará--Maranhão, Barreirinhas, Ceará e Potiguar. Essa região tem sido apontada como uma nova fronteira para a exploração offshore no país, reunindo condições geológicas promissoras e, simultaneamente, elevada sensibilidade ambiental.

Do ponto de vista físico, a Margem Equatorial apresenta grande variabilidade de lâmina d’água, sedimentos predominantemente moles e forte influência hidrodinâmica, fatores que afetam diretamente os mecanismos de propagação de vibração no solo e de ruído no meio aquático. Sob a perspectiva ecológica, trata-se de uma área de elevada biodiversidade, associada a rotas migratórias de grandes mamíferos marinhos e a ecossistemas ainda pouco impactados por atividades industriais intensivas.

Dessa forma, estudos voltados à modelagem da geração e propagação de vibração e ruído associados à instalação de estruturas offshore tornam-se particularmente relevantes, ao fornecer subsídios técnicos para a avaliação preventiva de impactos em regiões ambientalmente sensíveis, como a Margem Equatorial brasileira, e ao apoiar decisões de engenharia baseadas em critérios físicos e ambientais integrados.


% colocar algumas figuras

\begin{figure}
    \centering
    \includegraphics[width=0.9\linewidth]{Figura/margemeq.jpg}
    \caption{Localização da  Margem Equatorial.}
    \label{fig:equatorial}
\end{figure}
%https://nossaenergia.petrobras.com.br/w/nossas-atividades/margem-equatorial

\end{comment}